\section{Related work}

Graph anomaly detection covers anomalous nodes, edges, and subgraphs under structural, statistical, and learning-based formulations \cite{Akoglu_2014}. Temporal knowledge graph research has largely emphasized completion and forecasting, including logical-rule methods that provide paths or rules as explanations \cite{Liu_2022,Cai_2023}. These methods ask which future facts are likely to occur. Online anomaly detection has a different target: it must flag facts that should not be trusted while retaining an explanation useful for graph maintenance. AnoT is the closest base method because it learns a compact rule graph for online, interpretable anomaly scoring \cite{Zhang_2024}. \method keeps this rule-graph formulation rather than replacing it with a latent representation.

Rule miners such as AMIE+ search for compact, high-confidence symbolic implications in incomplete knowledge bases \cite{Galarraga_2015}. Our miner also returns a finite symbolic rule set, but its selection target is different: a temporal pattern must distinguish normal contexts from time-displaced references, then survive validation against evaluation-style time errors. Our temporal component is therefore not a generic contrastive representation learner. It uses contrast only to mine typed atomic-gap patterns and validates each candidate with explicit frequency, enrichment, posterior-direction, and MDL criteria. The result remains a rule file consumed by the same graph-based scorer. The workflow component is similarly conservative: its agents package, retrieve, verify, and explain artifacts. It does not delegate anomaly judgment to an unverified generative model.

The distinction from graph anomaly detection is also practical. A generic graph anomaly method often treats rarity, reconstruction error, or neighborhood inconsistency as the signal of interest \cite{Akoglu_2014}. In a temporal knowledge graph, an event can be rare but correct, or frequent but misplaced in time. CARGO-AnoT therefore does not use rarity as a sufficient explanation. Its score traces return a typed relation-role condition, a predecessor--target gap condition, or a context--target support condition. These conditions are narrower than a claim that an entire entity, relation, or subgraph is anomalous, which is appropriate for a curator deciding whether to inspect a single arriving quadruple.

Temporal link forecasting provides an adjacent but different form of temporal reasoning. Rule-based forecasters such as TLogic encode relation sequences that make a future link plausible \cite{Liu_2022}; survey work places these methods alongside embedding and neural approaches for temporal knowledge graph completion \cite{Cai_2023}. CARGO-AnoT reuses the useful part of this perspective---explicit temporal patterns---but changes the decision problem. A temporal-gap rule is not used to predict the most likely future object. It is used to assess whether a specified event occurs in a context and interval that are overrepresented among time-displaced references. The output remains an anomaly score with evidence, not a ranking over missing links.

Finally, the workflow differs from a free-form agent architecture. The term agentic describes ownership and traceability of intermediate artifacts, not autonomous problem reformulation. Candidate discovery is separable from evidence retrieval; evidence retrieval is separable from validation; and validation is separable from final scoring. This makes rejected rules informative rather than invisible. In the recorded time workflow, symbolic relation-pair and atomic-rule candidates were retained in the audit report even though none entered the final score. The distinction matters for research reporting: the paper can attribute the observed time gain to the validated contrastive atomic-gap source rather than suggesting that all conflict heuristics contributed.

